\section{Method}
\label{sec:method}

\begin{figure*}[htbp]
\centering
\includegraphics[width=\textwidth]{fig_overview.pdf}
\caption{The policy and its training loop. The shop state at one decision is
encoded as a set of tokens in five segments --- machines, jobs, vehicles, one
global token, and the aisle segments --- padded to a common width and passed
through a Transformer with no attention mask. One head scores the candidates of
each decision kind. Every decision of one episode is appended to a single
chain, which receives one terminal scalarized reward; the group-relative
advantage is computed over the chains sampled at that step, and the policy is
updated without a critic. The zone segment exists only when the congestion
mechanism is enabled; the other four are always present. The zone segment
together with the distance bias and the route head form the
transport-network-aware encoding of Section~\ref{sec:method:state}: they are
what puts the aisle network into the state rather than leaving it to be
inferred from the machine and vehicle tokens.}
\label{fig:overview}
\end{figure*}

\subsection{Overall framework}
\label{sec:method:framework}

Figure~\ref{fig:overview} shows the pipeline, and this subsection names its
stages in the order the running system passes through them.

\emph{Encoding.} At a decision instant the simulation hands over a snapshot:
machine states and buffers, job progress and due-date margins, vehicle
positions and loads, the shop clock, and the aisle occupancy. The snapshot is
turned into the token tensor of Section~\ref{sec:method:state}, with the
per-token features normalized by instance constants so that the same encoder
accepts a 10-job and a 20-job instance.

\emph{Scoring.} The encoder produces one embedding per token. The head of the
active decision kind scores its candidate set against those embeddings, and the
policy samples a candidate during training or takes the highest-scoring one
during evaluation.

\emph{Advancing.} The chosen candidate is applied to the simulation, which runs
until the next decision instant and then hands over a new snapshot. The chain
grows by one decision. The kind of the next decision is determined by the
triggering event, not by a fixed schedule, so the proportion of S, L, and
optional decisions in a chain is a property of the episode rather than a design
constant.

\emph{Learning.} The episode ends, one terminal reward is computed, and the
chains sampled at that step are compared with one another to give the advantage
of Algorithm~\ref{alg:train}.

\emph{Deployment.} At inference the policy needs the same snapshot, the same
candidate construction, and one forward pass; it does not need the reference
run, the reward, or the group. The candidate sets are built from the simulation
state exactly as during training, so the action space seen at deployment
matches the one seen during training. Evaluation in this paper uses argmax
action selection, which makes a trained policy deterministic given a seed; the
group-based sampling path exists only during training.

\subsection{Markov decision process formulation}
\label{sec:method:mdp}

We formulate the coupled problem as a Markov decision process
$\mathcal{M} = (\mathcal{S}, \mathcal{A}, P, r, \gamma)$ with a terminal
reward and no discounting.

\paragraph{State}
At a decision instant $t$ the state is the shop snapshot, encoded as a token
tensor and a segment assignment,
\begin{equation}
s_t = \left( \mathbf{X}_t \in \mathbb{R}^{N \times F_{\max}}, \;
            \mathbf{g}_t \in \{1,\dots,5\}^{N} \right),
\label{eq:state}
\end{equation}
where $N = m + n + V + 1 + Z$ is the number of tokens, $F_{\max}$ the padded
feature width, and $\mathbf{g}_t$ the segment identifier of each token. The
token set is built from the simulation state at the instant of the decision,
not from a fixed initial snapshot, so a decision taken late in an episode sees
the consequences of the decisions taken before it.
Section~\ref{sec:method:state} lists the features.

\paragraph{Action}
The action space is not a single set, because the decisions are of different
kinds. At each decision instant exactly one kind $k$ is active, drawn from
$\mathcal{K} = \{\mathrm{S}, \mathrm{L}, \mathrm{R}, \mathrm{M}, \mathrm{C},
\mathrm{B}\}$, and the action is a choice from a candidate set that depends on
the state and the kind:
\begin{equation}
a_t \in \mathcal{C}_k(s_t), \qquad
\mathcal{C}_k(s_t) = \{c_1, \dots, c_{C}\}, \quad C = C_k(s_t).
\label{eq:action}
\end{equation}
The candidate set varies in size: eligible machines for an operation, available
vehicles for a transport task, admissible paths for a travel, admissible batch
sizes for a pick-up. Formulating the action this way, rather than as a fixed
index into a problem-sized vector, is what lets one parameter set serve
instances of different size.

\paragraph{Transition}
$P(s_{t+1} \mid s_t, a_t)$ is the discrete-event simulation itself. The shop
advances until the next decision instant, and the triggering event --- an
operation released, a travel requested, a machine idle, a vehicle empty, a
batch starting --- determines the kind $k$ of the next decision. Because the
simulation is stochastic (failures, rework, repair durations), the transition
is stochastic even when the policy is deterministic.

\paragraph{Reward}
The reward is terminal and arrives once, at the end of the episode:
\begin{equation}
r(s_1, a_1, \dots, a_T) = w_1(-\Cmax) + w_2(-E) + w_3(-\TWT),
\label{eq:reward-mdp}
\end{equation}
with the weights of Equation~\eqref{eq:weights}. There is no intermediate
reward and no shaping term, so the credit for every decision in an episode is
assigned by the same scalar.

\paragraph{Objective}
The policy $\pi_\theta(a_t \mid s_t)$ is trained to maximize
$\mathbb{E}_{\pi_\theta}[r]$, with no discount factor, since the horizon is
the episode.

\subsection{Joint-chain formulation}
\label{sec:method:chain}

An episode of the scheduling problem contains many decisions of more than one
kind. Each operation needs a machine, and each transport task needs a vehicle.
A policy could be trained on each decision point separately, with one reward
per decision. We instead treat the whole episode as one decision chain.

Let $a_S$ be the sequence of machine choices, one per operation, and $a_L$ the
sequence of vehicle choices, one per transport task. A \emph{chain} is the pair
$(a_S, a_L)$ together with the optional decisions of
Section~\ref{sec:method:mechanisms}. Its log-probability under the policy is the
sum of the per-decision log-probabilities,
\begin{equation}
\begin{split}
\logp(\text{chain})
   &= \sum_{t} \logp_S(a_S[t] \mid s_t)
      + \sum_{t} \logp_L(a_L[t] \mid s_t) \\
   &\quad + \sum_{k} \sum_{t} \logp_k(a_k[t] \mid s_t),
\end{split}
\label{eq:chain-logp}
\end{equation}
where the last sum runs over the optional decision kinds $k$ that are enabled.
The chain receives a single terminal reward, computed after the episode ends.
No intermediate reward is used.

Three properties of the problem motivate this formulation.

First, group-relative advantages compare the members of a group with one
another. That comparison is only meaningful if every member was produced by the
same policy from the same state distribution. If the production and logistics
decisions were split between two policies, each group would mix samples drawn
from two different distributions, and the group mean would no longer estimate
a baseline for either of them. Multi-agent methods handle this by learning a
joint value function or a centralised critic~\citep{rashid2018qmix,yu2022mappo};
both routes reintroduce the critic that a group-relative method is meant to
avoid.

Second, the two decision layers are coupled by construction: the quality of a
dispatching decision depends on the production plan that generated the
transport tasks. A dispatching policy that is trained against a fixed or
randomly generated plan therefore learns against a task stream that the
deployed system will not produce.

Third, the coupling is asymmetric. Production decisions determine which
transport tasks exist and when they exist; transport decisions determine
whether the machines receive work on time. A chain-level advantage assigns
credit to the whole chain and leaves the split between the two layers to the
policy.

Treating a trajectory as one object and learning from it as a sequence is
itself an established formulation~\citep{chen2021decision}. What is specific
here is that the sequence spans two decision layers whose actions are of
different kinds and live in different candidate spaces, and that the reward
arrives only at the end. The candidate-scoring interface of
Section~\ref{sec:method:policy} is what lets one chain carry both.

\subsection{State representation}
\label{sec:method:state}

The policy reads the shop state as a sequence of tokens. Four kinds of entity
are encoded, plus one global token. Each token carries a fixed-width feature
vector. The feature vectors of the four segments are padded to a common width
$F_{\max}$ and concatenated into a single tensor of shape $(N, F_{\max})$, where
$N = m + n + V + 1$ for $m$ machines, $n$ jobs, and $V$ vehicles. A segment
identifier accompanies the tensor.

Machines, jobs, vehicles, and the global token have different numbers of
informative fields. A machine has 7, a job has 9, a vehicle has 11, and the
global token has 3. Because the four segments are padded to a common width,
column $c$ of a machine token and column $c$ of a job token can carry different
meanings. A learned type embedding, indexed by the segment identifier, is added
to every token and is what separates those meanings. Without it the shared
projection cannot distinguish a machine from a job.

Table~\ref{tab:features} lists the features. All quantities are normalized by
static instance constants --- total work content, the makespan of a reference
schedule, the machine count, the bounding box of the shop --- rather than by
per-episode statistics, because online decision making cannot observe the
statistics of an episode that has not finished.

\paragraph{Transport-network-aware encoding}
The four segments above describe the machines, the jobs, and the vehicles, and
none of them carries the shop's spatial structure: a machine token says how
much work is queued, not how far it is from anything else. We therefore add the
transport network to the state, in three parts that only work together. The
aisle segments become a fifth token segment, so each segment of the network is
an entity carrying its own occupancy and the time vehicles have waited on it. A
distance-dependent bias enters the attention score, so that a token can be
attended to in proportion to how far its entity is from the deciding one. And
the route head reads the segment tokens its candidate path would traverse,
which is what turns the first two parts from context into something the policy
can choose between. Section~\ref{sec:method:mechanisms} lists the switches that
enable each part.

This is a representation choice and not a claim about generalization. Whether
the encoding transfers to a shop layout it was not trained on is a separate
question, and this paper does not measure it.

\begin{table*}[htbp]
\centering
\caption{Token features. Each column is one scalar. Normalizers are instance
constants, not per-episode statistics.}
\label{tab:features}
\setlength{\tabcolsep}{4pt}
\begin{threeparttable}
\begin{tabular}{@{}llp{11.4cm}@{}}
\toprule
Segment & Width & Features \\
\midrule
Machine ($m$ tokens) & 7 &
  Queued processing work divided by total work content;
  queued operations divided by input capacity;
  output buffer occupancy divided by output capacity;
  busy indicator and remaining processing time as a fraction of the operation;
  remaining time until the next forced maintenance;
  failure rate divided by the largest failure rate in the instance. \\
Job ($n$ tokens) & 9 &
  Completed operations divided by total operations;
  remaining nominal processing time divided by total work content;
  due-date margin $(d_j - \text{now})/M_{\mathrm{ref}}$, which may be negative;
  completion indicator;
  current machine index, normalized, or $-1$ if not started;
  in-transit indicator and the index of the carrying vehicle;
  job weight divided by the largest weight. \\
Vehicle ($V$ tokens) & 11 &
  One-hot over idle, empty-travelling, loaded, and failed;
  current node coordinates divided by the bounding-box diagonal;
  assigned tasks divided by the instance maximum;
  battery fraction;
  load divided by capacity;
  speed multiplier divided by the largest multiplier;
  time spent waiting for an aisle segment divided by the wait limit. \\
Global (1 token) & 3 &
  Elapsed time divided by $M_{\mathrm{ref}}$;
  completed jobs divided by $n$;
  tasks in transit divided by the instance maximum. \\
\bottomrule
\end{tabular}
\begin{tablenotes}\footnotesize
\item The setup time of a job change is deliberately absent from the machine
      token. Setup is a property of a machine--job pair, not of either entity
      alone, so it is supplied to the machine-selection head as a candidate
      feature instead. A per-entity feature cannot express it.
\end{tablenotes}
\end{threeparttable}
\end{table*}

\subsection{Policy architecture}
\label{sec:method:policy}

The encoder is a standard Transformer~\citep{vaswani2017}. Tokens are
projected, given a type embedding, and passed through $L$ self-attention layers
with no mask, so every token can attend to every other token. We use $L = 8$
layers and model dimension $d = 128$ throughout.

Attention is unmasked rather than restricted to within-segment or
within-axis blocks. Segment identity is information about what a token is;
an attention mask is a restriction on what a token may see. The two are not
interchangeable, and dropping a restriction does not drop the information.

Each decision kind has its own head. A head scores \emph{candidates} rather
than a fixed action set, in the manner of pointer
networks~\citep{vinyals2015pointer}:
\begin{equation}
\ell = f_k\!\left(\mathbf{h}_{\text{token}}, \mathbf{g}_{\text{decision}},
   \mathbf{c}_{1:C}\right) \in \mathbb{R}^{C},
\label{eq:candidate}
\end{equation}
where $\mathbf{h}_{\text{token}}$ is the encoder output for the token that
owns the decision, $\mathbf{g}_{\text{decision}}$ collects decision-level
features, and $\mathbf{c}_{1:C}$ are per-candidate features. The number of
candidates $C$ varies with the state and with the instance, so the same
parameters serve a 10-job and a 20-job instance without resizing. Every head
reads the same token embeddings; no head receives a hand-built flat vector.

This matters because the candidate count is not a design constant. Eligible
machines per operation, available vehicles per task, viable paths per travel,
and admissible batch sizes all vary from decision to decision. Attention over
a set of elements is permutation-invariant by construction~\citep{lee2019set},
which is the property the head relies on: the score of a candidate must not
depend on where in the candidate list it happens to sit.

\subsection{Group-relative training over chains}
\label{sec:method:training}

For each training step the policy samples $G$ chains from the same state. Each
chain $g$ yields a terminal reward $r_g$. The advantage is the group-relative
score of Shao et al.~\citep{shao2024grpo},
\begin{equation}
A_g = \frac{r_g - \operatorname{mean}(r_{1:G})}
           {\operatorname{std}(r_{1:G}) + \epsilon},
\label{eq:advantage}
\end{equation}
and the objective is the REINFORCE objective with this advantage,
\begin{equation}
\mathcal{L} = -\frac{1}{G} \sum_{g=1}^{G} A_g \,
   \logp(\text{chain}_g),
\label{eq:loss}
\end{equation}
with $\logp$ from Equation~\eqref{eq:chain-logp}. No critic is used. The
optimizer is Adam. A single group is used per step ($J = 1$); the sampling
budget is spent on $G$ rather than on repeated seeds, because the group mean
already supplies the baseline and averaging over seeds would reduce noise
inside a quantity that is about to be compared within the group. The group is
also where diversity comes from: all $G$ chains start from the same state and
differ only in the sampling, the same device that multi-start policy
optimization uses to find several good solutions to one
instance~\citep{kwon2020pomo}.

\begin{algorithm}[t]
\caption{Joint-chain group-relative policy optimization}
\label{alg:train}
\begin{algorithmic}[1]
\Require instance, initial policy $\pi_\theta$, group size $G$, steps $T$
\Statex \textbf{reference} run the reference schedule once, cache
       $f^{\mathrm{ref}}$, $w$, $M_{\mathrm{ref}}$
\For{step $= 1$ \textbf{to} $T$}
  \For{$g = 1$ \textbf{to} $G$}
    \State sample a chain $\tau_g$: run the simulation under $\pi_\theta$
    \Statex \hspace{\algorithmicindent} and record every S, L, R, M, C, B decision
    \State $r_g \gets w_1(-\Cmax^{(g)}) + w_2(-E_g) + w_3(-\TWT_g)$
  \EndFor
  \State $\bar{r} \gets \frac{1}{G}\sum_{g} r_g$; \quad
         $\sigma \gets \sqrt{\frac{1}{G}\sum_{g} (r_g - \bar{r})^2}$
  \State $A_g \gets (r_g - \bar{r}) \,/\, (\sigma + \epsilon)$
         \Comment{advantage, Eq.~\eqref{eq:advantage}}
  \State $\mathcal{L} \gets -\frac{1}{G}\sum_{g} A_g \logp(\tau_g)$
  \State $\theta \gets \mathrm{Adam}\!\left(\theta, \nabla_\theta \mathcal{L}\right)$
\EndFor
\end{algorithmic}
\end{algorithm}

\paragraph{Training details}
The reference run is executed once per instance and configuration, and the
weights of Equation~\eqref{eq:weights} are derived from it before training
starts; the same run supplies the normalizing constant $M_{\mathrm{ref}}$ used
in the token features, so features and reward share one reference point.
Training uses $G = 8$ chains per step, $T = 300$ steps, Adam with a learning
rate of $3\times10^{-4}$, and a single pass over each group
($\mathrm{epochs} = 1$, so the default path is REINFORCE with a group-relative
advantage). Action sampling uses a seeded generator derived from the chain
index, which makes a run reproducible within one process but not across
processes. Evaluation uses argmax action selection over five perturbation
seeds, none of which is seen during training.

Equation~\eqref{eq:advantage} divides by the group standard deviation.
Section~\ref{sec:results} reports an ablation over this choice: a variant that
subtracts the group mean without dividing by the standard deviation, and a
variant that uses the raw reward. The three variants differ only in the
advantage; every switch and hyperparameter is otherwise identical. The same
comparison has been made for neural combinatorial optimization, where the
group-relative construction is used without a critic as
well~\citep{sepulveda2026baselinefree}, and the constrained variant of the same
family is what our budget mechanism
extends~\citep{hu2023constrained,achiam2017cpo,tessler2019rcpo}.

\paragraph{Reward}
The reward is a weighted scalarization of the three objectives,
\begin{equation}
r = w_1 (-\Cmax) + w_2 (-E) + w_3 (-\TWT),
\label{eq:reward}
\end{equation}
with weights derived from a reference schedule rather than chosen by hand:
\begin{equation}
w_i = \frac{1/f_i^{\mathrm{ref}}}{\sum_j 1/f_j^{\mathrm{ref}}},
\label{eq:weights}
\end{equation}
where $f_i^{\mathrm{ref}}$ is the measured value of objective $i$ under the
reference schedule. The reference schedule assigns every operation to its
fastest eligible machine and dispatches vehicles round-robin. The weights are
normalized so that improving each objective by one unit relative to the
reference contributes equally. We report the weight vector together with the
reference schedule, because the tardiness component of $f^{\mathrm{ref}}$
depends on the due-date calibration and the reward is not reproducible without
it.

Combining the objectives at the reward level keeps one scalar return per
episode, which is what the advantage estimator operates on; combining them at
the gradient level~\citep{sener2018mgda,yu2020pcgrad} would instead require
per-objective advantages and a rule for reconciling them.

The reward is terminal and undiscounted, and no shaping term is
added~\citep{ng1999shaping}. This is a deliberate simplification: a shaping
term would require a potential function over a state space in which the useful
potential is not obvious, and it would change what the episode-level comparison
inside a group means.

\subsection{Constraint mechanisms}
\label{sec:method:mechanisms}

Ten constraints are modelled. Making them switchable is not enough to make them
learnable; the policy must also be able to act on them. We group the mechanisms
by what they change.

\begin{itemize}
\item \textbf{Decision points.} The mechanism adds a new kind of decision that
      was previously made by a fixed rule, and the policy chooses it. The
      candidate-scoring interface of Equation~\eqref{eq:candidate} is what makes
      this a local change: a new head is added, the chain grows, and nothing
      else moves.
\item \textbf{Dynamics.} The mechanism changes how the shop evolves --- a
      breakdown, a returned task, a blocked buffer. The policy does not gain an
      action; it faces a different system.
\item \textbf{Input.} The mechanism adds features to the representation. No
      action and no trajectory changes.
\end{itemize}

The mechanisms are listed below. Every switch defaults to off,
and turning it off reproduces the pre-mechanism behaviour bit for bit; that
property is verified rather than asserted (Section~\ref{sec:setup:repro}).

\begin{description}\itemsep2pt
\item \textbf{Route choice (C1, decision point, \texttt{route\_k}).}
      At each travel the policy selects one of the $k$ shortest paths between
      the endpoints. Candidates carry path length, the number of aisle segments
      traversed, and the current contention on those segments.
\item \textbf{Zone tokens (C1, input, \texttt{route\_zones}).}
      Aisle segments become tokens, and the route head reads the segments its
      candidate traverses. Without this, the candidates of a travel differ only
      in their own summary features.
\item \textbf{Geometric bias (C1, input, \texttt{geom\_bias}).}
      A fixed, non-learned bias $-w\,d(i,j)/\text{diag}$ is added to the
      attention score, where $d$ is the shortest-path distance on the aisle
      graph. Nothing is learned here; the term only makes distance and
      congestion visible to attention.
\item \textbf{Maintenance head (C10, decision point, \texttt{pm\_head}).}
      When a machine is idle the policy chooses between starting maintenance
      now and deferring it. The forced-maintenance threshold remains a hard
      floor, so the policy can advance maintenance but not postpone it.
\item \textbf{Charging head (C9, decision point, \texttt{charge\_head}).}
      When a vehicle is idle the policy chooses between not charging and each
      charging station. A vehicle whose battery reaches zero becomes
      unavailable until it recovers.
\item \textbf{Multi-drop travel (C8, dynamics, \texttt{multi\_drop}).}
      One trip collects every queued task at a pick-up point and delivers to
      several drop points, bounded by capacity. Loaded travel becomes a
      sequence of segments rather than one.
\item \textbf{Batch head (C8, decision point, \texttt{batch\_head}).}
      At pick-up the policy chooses how many further tasks to append to the
      trip. Candidates carry batch size, the number of tasks whose relative
      order the batch would disturb, and the resulting travel saving.
\item \textbf{Failure failover (C7, dynamics, \texttt{agv\_failover}).}
      Tasks held by a failed vehicle are returned and offered to vehicles that
      are still running.
\item \textbf{Age-dependent failure (C3, dynamics, \texttt{machine\_age\_failure}).}
      The failure rate rises with the machine maintenance clock along a Weibull
      hazard, so the failure process acquires memory and maintenance acquires a
      second effect.
\end{description}

Four constraints need no dedicated mechanism. C2, C4, and C6 are already
visible in the token features of Table~\ref{tab:features}, and C5 is supplied
to the machine head as a candidate feature. Section~\ref{sec:results} reports
which of the mechanisms changed a measured quantity and which did not.

\paragraph{Why some ablation arms change more than one switch.}
A decision point that is enabled while the constraint it acts on is disabled
produces a dead action: the head samples a choice that cannot change anything,
and the chain log-probability still includes it. The advantage would then be
assigned partly to a decision with no effect, which contaminates the comparison
rather than merely wasting computation. Such combinations are therefore
rejected rather than permitted, and the consequence for the ablations is that
several arms necessarily change more than the one item they are named for. The
logistics arm, for instance, also disables the route head, the zone tokens, the
charging head, and the batch head, because each of those acts on a constraint
the arm has switched off. The released configuration files record the exact
switch set of every run, and Table~\ref{tab:ablation} marks the affected arms
so that the reader is not left to infer the switch set from the arm name.

\paragraph{Two mechanisms change the simulated system, not only the policy.}
Making vehicle capacity operational required changing the travel model: one
trip now collects every queued task at a pick-up point and delivers to several
drop points, so loaded travel is a sequence of segments rather than one. This
is a change to the shop dynamics and it is reported as such, because it means
the readings on either side of it are not comparable. The same is true of the
charging mechanism, which required that a vehicle at zero battery actually
become unavailable before there was a decision worth making. Both changes were
made because a parameter the model carried --- load capacity, battery state ---
had no effect on the simulation. The list below marks each
mechanism as a decision point, a dynamics change, or an input change for
exactly this reason.

\subsection{Lagrangian constraint budgets}
\label{sec:method:t3}

A constraint that is modelled but never violated contributes nothing to the
reward, so the policy has no reason to avoid violating it when the trade-off
turns favourable. We add a Lagrangian term with a per-constraint multiplier.
The construction follows the constrained policy optimization
lineage~\citep{achiam2017cpo,tessler2019rcpo}, with one difference that matters
here: those methods estimate the constraint cost with a value function, and we
do not use a critic. The multiplier is updated from the batch itself, which is
what makes the scheme compatible with a group-relative advantage.

Let $a_i$ be the activation of constraint $i$ over an episode, measured so
that it counts only what the policy was forced into: time spent waiting for an
aisle segment, minutes blocked by a full buffer, setup minutes, failures,
battery-depletion events, and maintenance triggered by the threshold. Let
$a_i^{\mathrm{ref}}$ be the same quantity under the reference schedule, so that
$\hat a_i = a_i / a_i^{\mathrm{ref}}$ is the activation as a fraction of the
reference level. With a budget $b_i$, the penalized reward of chain $g$ is
\begin{equation}
r'_g = r_g - \sum_i \lambda_i \hat a_{i,g},
\label{eq:lagrangian}
\end{equation}
and the advantage is computed from $r'$ as in Equation~\eqref{eq:advantage}.
The multipliers follow dual ascent,
\begin{equation}
\lambda_i \leftarrow \operatorname{clip}\!\left(
   \lambda_i + \eta \left(\hat a_i - b_i\right), \, 0, \, \lambda_{\max}\right),
\label{eq:dual}
\end{equation}
where $\hat a_i$ is the mean over the group, $\eta$ is a step size, and
$\lambda_{\max}$ guards against divergence. The multipliers persist across
training steps; they are not reset per step.

Two properties of this construction need to be stated because they bound what
it can do. The penalty is applied to the scalar reward \emph{before} the group
normalization, so the effective strength of $\lambda_i$ is scaled by the group
standard deviation of the reward and is not a direct statement of how much
weight the constraint carries. And the penalty enters the advantage only if the
group can express a difference: with $G = 2$ the normalized advantage is
identically $(-1,+1)$ regardless of the reward, so the penalty is inert. The
budget mechanism therefore requires $G \ge 3$.

Two constraints are excluded from the budget. Rework and vehicle failure are
memoryless events that the policy cannot influence, so penalizing them would
add noise to the reward rather than pressure to the policy. For the same
reason a constraint is admitted only when the policy has the corresponding
action available; the configuration is rejected at start-up if a constraint in
the budget has no enabled mechanism to act on it.

\paragraph{Calibration and the meaning of the budget.}
The reference activation $a_i^{\mathrm{ref}}$ is measured on the reference
schedule, and the budget $b_i$ is expressed as a fraction of it. The scale of
$a_i$ differs by three orders of magnitude across constraints --- aisle waiting
is measured in minutes, maintenance in events --- which is why the ratio form
is used rather than the raw quantity.

Calibration is done in a configuration in which all six constraints activate,
because on the standard parameters two of them never do
(Section~\ref{sec:results:active}) and a ratio with a zero denominator is
undefined. On the geometric caliber all ten instances yield six live
constraints. On the matrix caliber only four of the ten instances do, because
in the remaining six the buffer-blocking quantity is identically zero on the
reference schedule; those instances are excluded from the table rather than
entering it with an undefined reference, and a lookup on an absent instance
raises an error instead of falling back.

\paragraph{What the budget can and cannot reach.}
The penalty is computed from the activation of constraints in the
configuration being trained, not from the calibration configuration. If a
constraint never activates in training, its normalized activation is zero while
its reference is not, so its multiplier stays at zero and the constraint is not
penalized. Under the standard parameters this happens for two of the six
constraints, so the budget in practice acts on four of them. This is stated in
Section~\ref{sec:results:active} and is reflected in the way the budget arm is
described there.

\subsection{Computational cost}
\label{sec:method:cost}

The policy is evaluated once per decision, and a decision costs one forward
pass over the token sequence plus the construction of its candidate set. The
network dominates the cost of a step: across the configurations we measured,
the sampling pass, the recomputation pass, and the backward pass together
account for between 55 and 88 per cent of a step, and the simulation for the
remainder. Costs are reported per training step, each step comprising $G = 8$
sampled episodes and one update.

With the mechanisms switched off, a step costs between 0.85\,s and 6.00\,s
depending on instance size, and a step over all ten instances of the benchmark
costs 26.4\,s in total. The cost grows with the number of decisions per episode
rather than with the number of tokens, since each S and L decision is a
separate forward pass; the 20-job instances cost roughly five times the 10-job
instance for that reason.

The optional mechanisms add measured overhead, and one of them reduces cost.
Enabling the aisle-segment tokens multiplies the cost of a step by about 1.44,
because they lengthen the token sequence and the route head reads a subset of
them per candidate. The distance bias adds about 5\%, since the bias is
precomputed. Multi-drop travel and the batch head together add about 2\% on
MK01 and 7\% on MK10. Chunking the log-probability recomputation, by contrast,
reduced both memory and time on the largest instance: without it the peak
allocation reached 49.6\,GB and the run could not complete on the available
device, and with a chunk size of 128 the peak fell to 1.3\,GB and the step
became 1.36 times faster. All figures are single measurements on one machine
and are reported with their configurations in the released records; they should
be read as orders of magnitude rather than as benchmarks.
